% TOP_PROBLEM: {"id": 3000040, "title": "Gonality and edge subdivisions", "category_id": 2, "category": {"id": 2, "name": "combinatorics", "display_name": "Combinatorics", "description": "Counting problems, graph theory, discrete structures.", "slug": "combinatorics", "order_index": 2, "created_at": "2026-07-31T15:26:25.671Z"}, "status": "open", "problem_number": "AMR-029-0040", "set_id": 13, "statusQualification": "Uniform subdivision invariance is false; the definition retains the broader dependence question and the cited odd-segment question."}
% FIRST_POSTED: 2026-10-01
% ENTRY_KIND: statement_only
% UnsolvedMath ID 3000040; source catalogue code AMR-029-0040
% !TeX program = lualatex
% Definition and sources only; this file is not an LLM solution attempt.
\documentclass[11pt,a4paper]{article}
\usepackage[margin=25mm]{geometry}
\usepackage{fontspec}
\setmainfont{lmroman10-regular.otf}[BoldFont=lmroman10-bold.otf,
 ItalicFont=lmroman10-italic.otf,BoldItalicFont=lmroman10-bolditalic.otf]
\usepackage{amsmath,amssymb,mathtools}
\usepackage{unicode-math}
\setmathfont{latinmodern-math.otf}
\AtBeginDocument{\let\setminus\smallsetminus}
\usepackage{xurl,hyperref}
\hypersetup{colorlinks=true,urlcolor=blue}
\setcounter{secnumdepth}{0}
\setlength{\parindent}{0pt}
\setlength{\parskip}{5pt}
\setlength{\emergencystretch}{3em}
\begin{document}
\section*{Gonality and edge subdivisions}\label{3000040}
{\small UnsolvedMath ID 3000040 · AMR-029-0040 · Combinatorics · AMR Open Problem Lists}
\subsection{Definitions and mathematical statement}
For a finite connected loopless graph $G$, a divisor is an integer
vector $D\in\mathbb Z^{V(G)}$, of degree $\sum_vD(v)$. It is effective
if every coordinate is nonnegative. Two divisors are equivalent if
their difference is $L_Gz$ for an integer vector $z$, where
$(L_Gz)(v)=\sum_{vw\in E(G)}(z(v)-z(w))$ counts edge multiplicities.
An effective divisor has rank at least one if $D-\mathbf1_v$ is
equivalent to an effective divisor for every vertex $v$.
The divisorial gonality $\operatorname{gon}(G)$ is the minimum degree
of such a divisor.

Let $G^{(k)}$ replace every edge by a path of $k+1$ edges, inserting
$k$ new vertices per edge. Determine how $\operatorname{gon}(G^{(k)})$
depends on $G$ and $k$, including sharp bounds relative to
$\operatorname{gon}(G)$ and conditions for equality or strict decrease.
One specific remaining question asks whether strict decrease can
occur when $k+1>1$ is odd.

Uniform subdivision does not always preserve divisorial gonality:
the cited counterexamples already give a decrease when $k=1$.
Divisors on the subdivided graph may put chips on newly inserted
vertices, and equivalence uses its new Laplacian. This discrete problem
does not silently replace vertex-supported divisors by arbitrary
points of a metric graph.
\subsection{Short English statement}
How does subdividing every edge equally affect discrete divisorial
gonality, and can an odd number of replacement segments cause a decrease?
\subsection{Sources}
\begin{itemize}
\item[S1] ULAM.AI and the UnsolvedMath Contributors, supplied problems.json, record 3000040 (AMR-029-0040), AMR Open Problem Lists. Catalogue identity and source transcription; CC BY 4.0.
\url{https://huggingface.co/datasets/ulamai/UnsolvedMath}.
\item[S2] Egres Open - List of open problems, Open-problem page: Gonality and edge subdivisions. Original source indexed by AMR-029-0040.
\url{https://oldlemon.cs.elte.hu/egres/open/List_of_open_problems}.
\item[S3] EGRES Open, Gonality and edge subdivisions. Individual source page and explanatory remarks.
\url{https://lemon.cs.elte.hu/egres/open/Gonality_and_edge_subdivisions}.
\item[S4] J. van Dobben de Bruyn, H. Smit and M. van der Wegen, Discrete and metric divisorial gonality can be different, especially the final open questions.
\url{https://arxiv.org/abs/2106.12568}.
\end{itemize}
\end{document}
